\subsection{Complex Jordan Form and Block Sizes}

\begin{theorem} \label{thm:complex-jordan-form} \hlt{Jordan Form over $\mathbb{C}$}\\
Let $\dim_{\mathbb{C}}\mathcal{V}=n\geq1$ and let $T\in\LL(\mathcal{V})$. Then $\exists$ a basis $\mathcal{B}$ of $\mathcal{V}$ and a Jordan matrix $J$ with $[T]_{\mathcal{B}}=J$.
\begin{proof}
Let $\operatorname{spec}(T)=\{\lambda_1,\ldots,\lambda_m\}$, set $G_j=\mathcal{G}_{\lambda_j}(T)$, $N_j=(T-\lambda_jI_{\mathcal{V}})|_{G_j}$, and $T_j=T|_{G_j}$.\\
By Theorem \ref{thm:complex-primary-decomposition-generalized-eigenspaces}, $\mathcal{V}=\bigoplus\limits_{j=1}^mG_j$ and each $N_j$ is nilpotent. By Theorem \ref{thm:nilpotent-jordan-chain-basis}, choose an $N_j$-chain basis $\mathcal{B}_j$ of $G_j$.\\
For an $N_j$-chain $\mathcal{C}\subset\mathcal{B}_j$ of length $r$, $[N_j]_{\mathcal{C}}=J_r(0)$ and $[T_j]_{\mathcal{C}}=\lambda_jI_r+J_r(0)=J_r(\lambda_j)$. Let $\mathcal{B}=\mathcal{B}_1\sqcup\cdots\sqcup\mathcal{B}_m$.\\
By directness, $\mathcal{B}$ is a basis of $\mathcal{V}$, and $[T]_{\mathcal{B}}$ is a direct sum of Jordan blocks.\\
\hlt{Structural Commentary}: Primary decomposition reduces $T$ to nilpotent chain bases.\\
\hlt{Key Insight}: On $G_j$, $T=\lambda_jI+N_j$.
\end{proof}
\end{theorem}

\begin{theorem} \label{thm:real-jordan-block-form} \hlt{Real Jordan Block Form}\\
Let $n\in\mathbb{N}$, $A\in\operatorname{Mat}(n,n,\mathbb{R})$. For $a\in\mathbb{R}$, $b>0$, set $R(a,b)\defeq\left(\begin{smallmatrix}a&b\\-b&a\end{smallmatrix}\right)$.\\
Let $\mathcal{J}_k(a,b)=[C_{rs}]_{r,s=1}^k$ with $C_{rr}=R(a,b)$, $C_{r,r+1}=I_2$ for $r<k$, and $C_{rs}=0_2$ otherwise.\\
Then $\exists S\in\operatorname{GL}_n(\mathbb{R})$ such that $S^{-1}AS$ is a direct sum of blocks of the following types:
\begin{enumerate}[label=(\roman*)]
\setlength{\itemsep}{0pt}
\item $J_k(\lambda)$, where $k\in\mathbb{N}$ and $\lambda\in\mathbb{R}$;
\item $\mathcal{J}_k(a,b)$, where $k\in\mathbb{N}$, $a\in\mathbb{R}$, and $b>0$.
\end{enumerate}
\begin{proof} Complexify $A$ to $A_{\mathbb{C}}$. By Theorem \ref{thm:complex-jordan-form}, choose Jordan chains for $A_{\mathbb{C}}$.\\
By $\chi_A\in\mathbb{R}[t]$, $\lambda\notin\mathbb{R}$ implies $\overline\lambda\in\operatorname{spec}(A_{\mathbb{C}})$.
\begin{enumerate}[label=(\roman*)]
\setlength{\itemsep}{0pt}
\item For $\lambda\in\mathbb{R}$, let $G_{\lambda}^{\mathbb{R}}=\ker(A-\lambda I_n)^n$ and $N_{\lambda}=(A-\lambda I_n)|_{G_{\lambda}^{\mathbb{R}}}$. \\
Then $(G_{\lambda}^{\mathbb{R}})_{\mathbb{C}}=\ker(A_{\mathbb{C}}-\lambda I_n)^n$ and $N_{\lambda}$ is nilpotent.\\
By Theorem \ref{thm:nilpotent-jordan-chain-basis}, take an $N_{\lambda}$-chain basis $\mathcal{B}_{\lambda}$; $[A]_{\mathcal{B}_{\lambda}}$ is a direct sum of $J_k(\lambda)$ blocks.
\item Let $\lambda=a+ib$, $b>0$, and let $v=(v_1,\ldots,v_k)$ be a $\lambda$-chain. By Lemma \ref{lem:complexification-conjugate-generalized-eigenspaces}, $\overline v$ is a $\overline\lambda$-chain.\\
Let $v_j=u_j+iw_j$; $U_j=[u_j\ w_j]$, $U_0=0$, $\mathcal{U}_v=(u_1,w_1,\ldots,u_k,w_k)$,  $W_v=\operatorname{span}_{\mathbb{C}}(v,\overline v)\cap\mathbb{R}^n$.\\
By $(v,\overline v)\leftrightarrow\mathcal{U}_v$, $\mathcal{U}_v$ is real basis of $W_v$. By chain relations, $AU_j=U_jR(a,b)+U_{j-1}$; $[A]_{\mathcal{U}_v}=\mathcal{J}_k(a,b)$.
\end{enumerate}
Let $\mathcal{B}$ concatenate all $\mathcal{B}_{\lambda}$ and one $\mathcal{U}_v$ for each pair $(\lambda,\overline\lambda)$, $\operatorname{Im}\lambda>0$.\\
By primary directness, $\mathcal{B}$ is a real basis of $\mathbb{R}^n$. With $S=[\mathcal{B}]$, $S^{-1}AS$ has the asserted blocks.\\
\hlt{Structural Commentary}: Real chains and conjugate-pair chains exhaust the primary decomposition.\\
\hlt{Key Insight}: A nonreal chain becomes one real block chain.
\end{proof}
\end{theorem}

\begin{figure}[H]
  \centering
  \begin{tikzpicture}[
  font=\small,
  block/.style={draw=blue!65!black, fill=blue!6, rounded corners=2pt, minimum width=2.15cm, minimum height=0.68cm, align=center},
  arrow/.style={-{Stealth[length=2mm]}, thick, gray!75}
]
  \node[block] (v) at (0,0) {$\mathcal{V}$};
  \node[block] (g1) at (3.0,0.85) {$\mathcal{G}_{\lambda_1}(T)$};
  \node[block] (gj) at (3.0,0) {$\cdots$};
  \node[block] (gm) at (3.0,-0.85) {$\mathcal{G}_{\lambda_m}(T)$};

  \node[block] (j1) at (6.0,0.85) {$J_k(\lambda_1)$};
  \node[block] (jj) at (6.0,0) {$\cdots$};
  \node[block] (jm) at (6.0,-0.85) {$J_k(\lambda_m)$};

  \draw[arrow] (v.east) -- (g1.west);
  \draw[arrow] (v.east) -- (gj.west);
  \draw[arrow] (v.east) -- (gm.west);
  \draw[arrow] (g1.east) -- (j1.west);
  \draw[arrow] (gj.east) -- (jj.west);
  \draw[arrow] (gm.east) -- (jm.west);
\end{tikzpicture}

  \caption{Primary pieces $\mathcal{G}_{\lambda_j}(T)$ give nilpotent chains, hence Jordan blocks.}
  \label{fig:jordan-form-primary-chains}
\end{figure}

\begin{corollary} \label{cor:jordan-blocks-diagonalizability} \hlt{Diagonalizability in Jordan Form}\\
Let $n\in\mathbb{N}$, $A\in\operatorname{Mat}(n,n,\mathbb{C})$, $J=\bigoplus\limits_{\alpha=1}^s J_{k_{\alpha}}(\mu_{\alpha})$ be Jordan form of $A$.\\
Set $\delta_{\lambda\alpha}\defeq\delta_{\lambda,\mu_{\alpha}}$ and $b_J(\lambda)=\sum\limits_{\alpha=1}^s\delta_{\lambda\alpha}$ for $\lambda\in\operatorname{spec}(A)$. Then:
\begin{enumerate}[label=(\roman*)]
\setlength{\itemsep}{0pt}
\item $\dim E_{\lambda}(A)=b_J(\lambda)$, $\lambda\in\operatorname{spec}(A)$;
\item $A$ is diagonalizable if and only if every Jordan block in $J$ has size $1$.
\end{enumerate}
\begin{proof}
Choose $S\in\operatorname{GL}_n(\mathbb{C})$ such that $A=SJS^{-1}$.
\begin{enumerate}[label=(\roman*)]
\setlength{\itemsep}{0pt}
\item For $\lambda\in\operatorname{spec}(A)$, $S:E_{\lambda}(J)\cong E_{\lambda}(A)$ and $\dim E_{\lambda}(J_{k_{\alpha}}(\mu_{\alpha}))=\delta_{\lambda\alpha}$. Thus $\dim E_{\lambda}(A)=\sum\limits_{\alpha=1}^s\delta_{\lambda\alpha}=b_J(\lambda)$.
\item Put $a_J(\lambda)=\sum\limits_{\alpha=1}^s k_{\alpha}\delta_{\lambda\alpha}$. By Theorem \ref{thm:diagonalizable-iff-algebraic-geometric-multiplicities-match} and (i), $A$ is diagonalizable $\Longleftrightarrow a_J=b_J$ on $\operatorname{spec}(A)$.\\
Since $a_J(\lambda)-b_J(\lambda)=\sum\limits_{\alpha=1}^s(k_{\alpha}-1)\delta_{\lambda\alpha}$ and $k_{\alpha}\geq1$, this is equivalent to $k_{\alpha}=1$ for every $\alpha$.
\end{enumerate}
\hlt{Structural Commentary}: Similarity transfers eigenspaces; block kernels count directions.\\
\hlt{Key Insight}: Algebraic multiplicity counts block lengths.
\end{proof}
\end{corollary}

\begin{lemma} \label{lem:nilpotent-jordan-block-power-rank} \hlt{Powers of a Nilpotent Jordan Block}\\
Let $k\in\mathbb{N}$, let $p\in\mathbb{N}_0$, and set $N=J_k(0)\in\operatorname{Mat}(k,k,\mathbb{C})$. Then:
\begin{enumerate}[label=(\roman*)]
\setlength{\itemsep}{0pt}
\item $\operatorname{rank}(N^p)=\max(k-p,0)$;
\item $\dim\ker(N^p)=\min(p,k)$.
\end{enumerate}
\begin{proof}
Let $(\mathbf{e}_j)_{j=1}^k$ be the standard basis. Set $(q)_+\defeq\max(q,0)$ and $\mathcal{E}_m=\operatorname{span}(\mathbf{e}_1,\ldots,\mathbf{e}_m)$ for $0\leq m\leq k$, with $\mathcal{E}_0=\{\mathbf{0}\}$. By $N\mathbf{e}_j=\mathbf{e}_{j-1}$ for $j>1$ and $N\mathbf{e}_1=\mathbf{0}$, $N^p\mathbf{e}_j=\mathbf{e}_{j-p}$ for $j>p$ and $\mathbf{0}$ otherwise.
\begin{enumerate}[label=(\roman*)]
\setlength{\itemsep}{0pt}
\item $\operatorname{im}(N^p)=\mathcal{E}_{(k-p)_+}$, hence $\operatorname{rank}(N^p)=\dim\mathcal{E}_{(k-p)_+}=(k-p)_+=\max(k-p,0)$.
\item By rank-nullity and (i), $\dim\ker(N^p)=k-(k-p)_+=\min(p,k)$.
\end{enumerate}
\hlt{Structural Commentary}: Powers shift the chain and truncate its span.\\
\hlt{Key Insight}: Rank counts surviving positions.
\end{proof}
\end{lemma}

\begin{definition} \label{def:weyr-characteristic} \hlt{Weyr Characteristic}\\
Let $A\in\operatorname{Mat}(n,n,\mathbb{C})$ and let $\lambda\in\operatorname{spec}(A)$.
\begin{enumerate}[label=(\roman*)]
\setlength{\itemsep}{0pt}
\item For $p\in\mathbb{N}_0$, define the shifted-power rank $r_p(\lambda,A)\defeq\operatorname{rank}\bigl((A-\lambda I_n)^p\bigr)$.
\item For $p\in\mathbb{N}$, define the Weyr number $w_p(\lambda,A)\defeq r_{p-1}(\lambda,A)-r_p(\lambda,A)$. If $s\defeq\max\{p\in\mathbb{N}:w_p(\lambda,A)>0\}$, then the finite sequence $\bigl(w_1(\lambda,A),\ldots,w_s(\lambda,A)\bigr)$ is the \hlt{Weyr characteristic} of $A$ associated with $\lambda$.
\end{enumerate}
\end{definition}

\begin{theorem} \label{thm:jordan-block-sizes-from-rank-drops} \hlt{Jordan Block Sizes from Rank Drops}\\
Let $n\in\mathbb{N}$, $A\in\operatorname{Mat}(n,n,\mathbb{C})$, $\lambda\in\operatorname{spec}(A)$, and $J=\bigoplus\limits_{\alpha=1}^sJ_{k_{\alpha}}(\mu_{\alpha})$ be a Jordan form of $A$.\\
For $q\in\mathbb{N}_0$, set $r_q=\operatorname{rank}\bigl((A-\lambda I_n)^q\bigr)$; for $q\in\mathbb{N}$, set $w_q=r_{q-1}-r_q$.\\
For $p\in\mathbb{N}$, set $I_p=\{\alpha:\mu_{\alpha}=\lambda,\ k_{\alpha}\geq p\}$. Then:
\begin{enumerate}[label=(\roman*)]
\setlength{\itemsep}{0pt}
\item $w_p=\# I_p$;
\item $w_p-w_{p+1}=\#(I_p\setminus I_{p+1})=\#\{\alpha:\mu_{\alpha}=\lambda,\ k_{\alpha}=p\}$.
\end{enumerate}
\begin{proof}
By similarity-invariance of rank, assume $A=J$. Let $\mathbf{1}_E$ be the indicator of $E$.\\
Set $B_{\alpha}=J_{k_{\alpha}}(\mu_{\alpha})$ and $\rho_q(B)=\operatorname{rank}\bigl((B-\lambda I)^q\bigr)$.\\
Put $\Delta_{p,\alpha}=\rho_{p-1}(B_{\alpha})-\rho_p(B_{\alpha})$. By rank additivity, $w_p=\sum\limits_{\alpha=1}^s\Delta_{p,\alpha}$.
\begin{enumerate}[label=(\roman*)]
\setlength{\itemsep}{0pt}
\item By invertibility for $\mu_{\alpha}\neq\lambda$ and by Lemma \ref{lem:nilpotent-jordan-block-power-rank} for $\mu_{\alpha}=\lambda$, $\Delta_{p,\alpha}=\mathbf{1}_{I_p}(\alpha)$. Hence $w_p=\sum\limits_{\alpha=1}^s\mathbf{1}_{I_p}(\alpha)=\#I_p$.
\item By (i), $w_p-w_{p+1}=\#I_p-\#I_{p+1}=\#(I_p\setminus I_{p+1})=\#\{\alpha:\mu_{\alpha}=\lambda,\ k_{\alpha}=p\}$.
\end{enumerate}
\hlt{Structural Commentary}: Rank drops split into blockwise indicators.\\
\hlt{Key Insight}: A block contributes until its length.
\end{proof}
\end{theorem}

\begin{figure}[H]
  \centering
  \begin{tikzpicture}[
  font=\small,
  rank/.style={draw=blue!60!black, fill=blue!6, rounded corners=2pt, minimum width=0.96cm, minimum height=0.50cm, align=center},
  count/.style={draw=red!60!black, fill=red!6, rounded corners=2pt, minimum width=2.90cm, minimum height=0.50cm, align=center},
  arrow/.style={-{Stealth[length=2mm]}, thick, gray!75}
]
  \node[rank] (r0) at (0,0) {$r_0$};
  \node[rank] (r1) at (1.36,0) {$r_1$};
  \node[rank] (r2) at (2.72,0) {$r_2$};
  \node[rank] (dots) at (4.08,0) {$\cdots$};
  \node[rank] (rq) at (5.44,0) {$r_q$};

  \node[count] (drop) at (2.72,-0.84) {$w_p=r_{p-1}-r_p$};
  \node[count, minimum width=3.28cm] (exact) at (2.72,-1.55) {size-$p$ blocks $=w_p-w_{p+1}$};

  \draw[arrow] (r0) -- (r1);
  \draw[arrow] (r1) -- (r2);
  \draw[arrow] (r2) -- (dots);
  \draw[arrow] (dots) -- (rq);
  \draw[arrow] (drop) -- (exact);
\end{tikzpicture}

  \caption{For $r_p=\operatorname{rank}((A-\lambda I)^p)$, drops count surviving chains; $w_p-w_{p+1}$ counts size-$p$ blocks.}
  \label{fig:weyr-rank-drops-jordan-blocks}
\end{figure}

\begin{corollary} \label{cor:jordan-form-unique-up-to-block-permutation} \hlt{Uniqueness of Jordan Form}\\
Let $n\in\mathbb{N}$, $A\in\operatorname{Mat}(n,n,\mathbb{C})$. Let $J,J'$ be Jordan matrices similar to $A$.\\
For $M\in\{J,J'\}$, let $\mathfrak{B}(M)$ be its multiset of Jordan blocks, set $m_M(\lambda,k)=\#\{B\in\mathfrak{B}(M):B=J_k(\lambda)\}$. Then:
\begin{enumerate}[label=(\roman*)]
\setlength{\itemsep}{0pt}
\item $m_J(\lambda,k)=m_{J'}(\lambda,k)$ for all $(\lambda,k)\in\operatorname{spec}(A)\times\mathbb{N}$;
\item $\exists$ a block-permutation matrix $P$ such that $J'=P^{-1}JP$.
\end{enumerate}
\begin{proof}
Set $d_k^{\lambda}(M)=w_k(\lambda,M)-w_{k+1}(\lambda,M)$. By similarity invariance, $d_k^{\lambda}(J)=d_k^{\lambda}(J')$ $\forall (\lambda,k)\in\operatorname{spec}(A)\times\mathbb{N}$.
\begin{enumerate}[label=(\roman*)]
\setlength{\itemsep}{0pt}
\item For $(\lambda,k)\in\operatorname{spec}(A)\times\mathbb{N}$, by Theorem \ref{thm:jordan-block-sizes-from-rank-drops}, $m_J(\lambda,k)=d_k^{\lambda}(J)=d_k^{\lambda}(J')=m_{J'}(\lambda,k)$.
\item By (i), $\mathfrak{B}(J)=\mathfrak{B}(J')$, so $\exists$ a block-permutation matrix $P$ with $J'=P^{-1}JP$.
\end{enumerate}
\hlt{Structural Commentary}: Shifted ranks determine block multiplicities.\\
\hlt{Key Insight}: Weyr data fixes Jordan blocks.
\end{proof}
\end{corollary}

\begin{method} \label{meth:jordan-form-from-rank-drops} \hlt{Computing Jordan Blocks by Rank Drops}\\
Let $A\in\operatorname{Mat}(n,n,\mathbb{C})$ and let $\lambda\in\operatorname{spec}(A)$. To determine the $\lambda$-Jordan blocks:
\begin{enumerate}[label=(\roman*)]
\setlength{\itemsep}{0pt}
\item Compute $r_p(\lambda,A)=\operatorname{rank}\bigl((A-\lambda I_n)^p\bigr)$ for $p=0,1,2,\ldots$ until $r_{p-1}(\lambda,A)=r_p(\lambda,A)$.
\item Set $w_p(\lambda,A)=r_{p-1}(\lambda,A)-r_p(\lambda,A)$. This is the number of $\lambda$-blocks of size at least $p$.
\item For $p\in\mathbb{N}$, take $w_p(\lambda,A)-w_{p+1}(\lambda,A)$ copies of $J_p(\lambda)$.
\item Repeat (i)--(iii) as $\lambda$ ranges over $\operatorname{spec}(A)$, and take the direct sum of all resulting blocks.
\end{enumerate}
The validity of steps (ii)--(iii) follows from Theorem \ref{thm:jordan-block-sizes-from-rank-drops}.
\end{method}
